%\ifodd\value{page}\hbox{}\newpage\fi
\chapter*{Dhyanam 2}

\begin{sloka}
{\sanskritfont
\begin{tabular}{c}
अरुणां करुणातरङ्गिताक्षीं \\
धृतपाशाङ्कुशपुष्पबाणचापाम् ।\\
अणिमादिभिरावृतां मयूखैः \\
अहमित्येव विभावये भवानीम् ॥
\end{tabular}

\vspace{1.5em}

{\small \iastfont
aruṇāṃ karuṇā-taraṅgitākṣīṃ \\
dhṛta-pāśāṅkuśa-puṣpa-bāṇa-cāpām |\\
aṇimādibhir āvṛtāṃ mayūkhaiḥ \\
aham ity eva vibhāvaye bhavānīm ||\\
}

\vspace{1em}

{\small 
\anvaya{%
अरुणाम् करुणातरङ्गिताक्षीम् पाशाङ्कुशपुष्पबाणचापधृताम्
अणिमादिभिः मयूखैः आवृताम् भवानीम् अहम् इति एव विभावये ॥
}
}

\small \iastfont \itmean{%
«Contemplo Bhavānī come rossa, dagli occhi increspati di compassione,
che regge il laccio, il pungolo, l’arco e la freccia fiorita;
avvolta dai raggi dei poteri sottili come aṇimā,
e la riconosco interiormente proprio come “io stesso”.»%
}


}
\end{sloka}

\wordmeaning{%
\wordline{अरुणाम्}{aruṇām}
  { rossa, color dell’aurora;}%
\wordline{करुणा-तरङ्गित-अक्षीम्}{karuṇā-taraṅgitākṣīm}
  { dagli occhi increspati di compassione;}%
\wordline{धृत}{dhṛta}
  { che regge, che tiene;}%
\wordline{पाश}{pāśa}
  { il laccio;}%
\wordline{अङ्कुश}{aṅkuśa}
  { il pungolo;}%
\wordline{पुष्प-बाण}{puṣpa-bāṇa}
  { la freccia fiorita;}%
\wordline{चाप}{cāpa}
  { l’arco;}%
\wordline{अणिमादिभिः}{aṇimādibhiḥ}
  { dai poteri sottili come aṇimā e gli altri;}%
\wordline{मयूखैः}{mayūkhaiḥ}
  { da raggi, irradiazioni;}%
\wordline{आवृताम्}{āvṛtām}
  { avvolta, circondata;}%
\wordline{अहम् इति एव}{aham iti eva}
  { proprio come “io stesso”;}%
\wordline{विभावये}{vibhāvaye}
  { riconosco interiormente, contemplo;}%
\wordline{भवानीम्}{bhavānīm}
  { Bhavānī, la Dea Suprema;}%
}

